% This file was converted to LaTeX by Writer2LaTeX ver. 1.9.9
% see http://writer2latex.sourceforge.net for more info
\documentclass{article}
\usepackage{calc,amsmath,amssymb,amsfonts}
\usepackage[T1]{fontenc}
\usepackage[italian]{babel}
\usepackage[style=numeric,backend=biber]{biblatex}
\author{utente}
\date{2026-08-18}
\begin{document}
Capitolo 2 – La prima automobile

Per oltre quarant'anni il videogioco che avevo immaginato da ragazzo non aveva avuto una sola riga di codice. Era
rimasto confinato nelle conversazioni con gli amici, nelle idee appuntate soltanto nella memoria e in quella
convinzione, tipica di quando si è giovani, che prima o poi sarebbe arrivato il momento giusto per cominciare. Quel
momento non era arrivato allora, ma adesso il Commodore 64 era nuovamente davanti a me e non avevo più alcuna ragione
per rimandare.

Restava però una domanda che, per un programmatore, è spesso più difficile di quanto sembri: da dove si comincia?

Se pensiamo al gioco nella forma che vorremmo raggiungere, la risposta rischia di diventare immediatamente complicata.
Immaginiamo una strada che scorre, l'automobile del giocatore, i veicoli avversari, gli ostacoli ai lati della
carreggiata, gli spari, il punteggio e, più avanti, suoni ed effetti capaci di dare al tutto il ritmo di un vero gioco
d'azione. Sono elementi che nella nostra mente esistono già contemporaneamente, perché possiamo immaginare il risultato
finale molto prima di essere in grado di costruirlo.

Un programma, invece, ci costringe a procedere in un altro modo. Prima della strada deve esistere qualcosa che possa
percorrerla; prima degli avversari deve esserci il veicolo del giocatore; prima di leggere il joystick dobbiamo avere
qualcosa da muovere. La nostra prima milestone, che chiameremo semplicemente M01, avrà quindi un obiettivo molto più
modesto dell'intero videogioco: fare comparire sullo schermo la nostra automobile.

Potrebbe sembrare poco. In realtà sarà sufficiente per costringerci a entrare per la prima volta dentro il funzionamento
del VIC-II e a scoprire uno degli strumenti che hanno reso il Commodore 64 una macchina particolarmente adatta ai
videogiochi.

Prima prepariamo lo schermo

Il nostro primo listato comincia senza occuparsi ancora dell'automobile. Prima di mostrare qualsiasi cosa, vogliamo
infatti sapere in quale condizione si trova lo schermo sul quale lavoreremo.

Le prime istruzioni operative della M01 sono queste:

60 rem pulisce lo schermo

70 print chr\$(147)

80 rem bordo nero

90 poke 53280,0

100 rem sfondo grigio chiaro

110 poke 53281,15

PRINT CHR\$(147) invia al Commodore il carattere di controllo che pulisce lo schermo. Potremmo eseguire il programma
anche senza questa istruzione, ma il risultato dipenderebbe da ciò che era presente sul monitor prima del RUN.
Preferiamo invece che sia il programma stesso a costruire, per quanto possibile, il proprio stato iniziale.

Le due istruzioni successive ci fanno incontrare una parola che utilizzeremo moltissime volte: POKE.

Nel BASIC del Commodore 64, POKE permette di scrivere direttamente un valore in una determinata locazione di memoria. La
sua forma è molto semplice:

POKE indirizzo,valore

Dietro questa semplicità si nasconde però una delle caratteristiche fondamentali della macchina. Sul Commodore 64 alcune
locazioni di memoria non servono soltanto a conservare numeri: sono collegate direttamente ai dispositivi hardware.
Scrivere un valore in uno di questi indirizzi significa quindi modificare il comportamento di un componente del
computer.

L'indirizzo 53280, per esempio, corrisponde al registro che controlla il colore del bordo dello schermo. Scrivendovi il
valore zero scegliamo il nero:

90 poke 53280,0

La locazione successiva, 53281, controlla invece il colore dello sfondo. Il valore 15 rappresenta il grigio chiaro:

110 poke 53281,15

Non abbiamo ancora disegnato nulla, ma abbiamo già fatto qualcosa di importante: per la prima volta il nostro programma
non si è limitato a stampare un testo attraverso il BASIC. Ha scritto direttamente in due registri del chip video e il
VIC-II ha reagito modificando ciò che vediamo.

Molti anni fa Rossi Brunori aveva usato parole molto più semplici per farmi comprendere questo principio. Il computer,
diceva, fa quello che gli scriviamo. Allora lo avevo visto con un PRINT {\textquotedbl}CIAO{\textquotedbl}. Adesso
stiamo cominciando a scoprire che alcune delle conversazioni più interessanti con il Commodore non passano
necessariamente dalle parole del BASIC: a volte basta scrivere il numero giusto nel posto giusto.

Un'automobile fatta di numeri

Preparato lo schermo, possiamo finalmente occuparci del protagonista della nostra prima milestone.

Il VIC-II mette a disposizione otto oggetti grafici indipendenti dallo schermo testuale: gli sprite. Per il momento
possiamo immaginarli come piccoli disegni che il chip video è in grado di sovrapporre al resto dell'immagine e spostare
modificandone semplicemente le coordinate.

È esattamente ciò che ci serve. La nostra automobile dovrà muoversi continuamente quando introdurremo il joystick e, più
avanti, dovrà convivere con veicoli avversari e altri elementi del gioco. Disegnarla direttamente dentro lo schermo
sarebbe possibile, ma significherebbe dover cancellare e ridisegnare continuamente la sua immagine. Uno sprite nasce
invece proprio per rappresentare oggetti di questo tipo.

Uno sprite standard del Commodore 64 misura 24 punti in larghezza e 21 in altezza. In modalità monocromatica ogni punto
corrisponde a un singolo bit: acceso o spento. Poiché una riga contiene 24 bit e un byte ne contiene 8, per descrivere
una riga sono necessari tre byte:

24 / 8 = 3 byte

Moltiplicandoli per le 21 righe otteniamo:

21 × 3 = 63 byte

Ed ecco spiegato il numero che compare nel primo ciclo importante della M01:

130 rem carica i 63 byte dello sprite

140 rem a partire da \$0340 = 832

150 for i=0 to 62

160 read a

170 poke 832+i,a

180 next i

Il ciclo parte da zero e arriva a 62. È facile, soprattutto le prime volte, guardare quei due valori e pensare che siano
soltanto 62 iterazioni; contando anche lo zero, invece, ne otteniamo esattamente 63.

Ad ogni passaggio READ A legge uno dei valori memorizzati nelle istruzioni DATA che troveremo alla fine del programma.
La riga successiva lo copia in memoria:

170 poke 832+i,a

Alla prima iterazione I vale zero e quindi il primo valore viene scritto all'indirizzo 832. Alla seconda I vale 1 e
scriviamo nell'833. Il procedimento continua fino a quando l'ultimo dei 63 byte viene depositato all'indirizzo 894.

È interessante osservare che le istruzioni DATA possono trovarsi anche molto più avanti nel listato. READ non pretende
che siano collocate immediatamente dopo il ciclo: il BASIC le cerca nel programma e ne preleva i valori nell'ordine in
cui sono stati scritti. Questo ci permette di tenere separata la parte esecutiva dalla lunga sequenza numerica che
descrive graficamente l'automobile.

Resta però una domanda: perché abbiamo scelto proprio l'indirizzo 832?

Il blocco numero 13

Gli sprite non possono essere memorizzati in un punto arbitrario della memoria. Il VIC-II li cerca all'inizio di blocchi
da 64 byte. Il nostro disegno ne utilizza 63; il blocco che lo contiene, però, ne occupa comunque 64.

L'indirizzo 832 è particolarmente comodo perché è esattamente divisibile per 64:

832 / 64 = 13

Il nostro sprite occupa quindi il blocco numero 13.

Questa relazione viene riportata esplicitamente nel programma:

200 rem blocco sprite: 832/64 = 13

210 rem puntatore sprite 0 = 2040

220 poke 2040,13

Abbiamo copiato l'immagine a partire dall'indirizzo 832, ma il VIC-II non può sapere autonomamente che proprio lì
abbiamo deciso di collocare l'automobile. Dobbiamo indicarglielo attraverso un puntatore sprite.

Nella configurazione standard con cui stiamo lavorando, il puntatore dello sprite 0 si trova all'indirizzo 2040.
Inserendo il valore 13:

220 poke 2040,13

stiamo dicendo al VIC-II che i dati grafici del primo sprite devono essere cercati nel blocco numero 13.

La scelta del termine puntatore comincia quindi a diventare intuitiva. Quel byte non contiene l'immagine e non contiene
nemmeno l'indirizzo 832 scritto per esteso; contiene un numero che permette al VIC-II di arrivare al blocco nel quale
si trova l'immagine.

È anche il primo esempio di una caratteristica che incontreremo continuamente lavorando a basso livello sul Commodore
64: risparmiare memoria significa spesso rappresentare un'informazione nel modo più compatto possibile. Se tutti gli
sprite devono iniziare su confini da 64 byte, non occorre memorizzarne l'indirizzo completo. Basta indicare quale
blocco utilizzare.

Dove deve apparire?

A questo punto la nostra automobile esiste in memoria e il VIC-II sa dove recuperarne l'immagine. Non abbiamo ancora
stabilito, però, dove debba comparire.

Lo sprite 0 possiede una coordinata orizzontale X e una verticale Y. Per questa prima milestone scegliamo semplicemente
una posizione che ci permetta di osservare comodamente l'automobile:

240 rem posizione iniziale

250 poke 53248,172

260 poke 53249,120

Il registro 53248 contiene la coordinata X dello sprite 0, mentre 53249 ne contiene la coordinata Y. Impostiamo quindi X
a 172 e Y a 120.

Questi valori non hanno ancora nulla a che fare con la futura carreggiata. Non stiamo decidendo quale sarà la posizione
di partenza definitiva del giocatore: stiamo semplicemente cercando un punto adatto a verificare che il nostro primo
sprite venga visualizzato correttamente.

La coordinata X introduce però una piccola particolarità.

Un byte può contenere valori compresi tra 0 e 255, mentre la coordinata orizzontale degli sprite del VIC-II può superare
255. Per rappresentare quel valore aggiuntivo esiste quindi un altro registro, all'indirizzo 53264, nel quale ciascun
bit fornisce il nono bit della coordinata X di uno degli otto sprite.

Nella M01 la nostra X vale soltanto 172 e non abbiamo bisogno di quel bit aggiuntivo. Lo azzeriamo esplicitamente:

280 rem azzera il nono bit della x

290 poke 53264,0

Questa istruzione merita una piccola annotazione in prospettiva. Scrivere zero nell'intero registro significa azzerare
il nono bit della coordinata X di tutti gli sprite, non soltanto dello sprite 0. Adesso non rappresenta un problema,
perché esiste soltanto la nostra automobile. Quando introdurremo altri sprite non potremo più permetterci di modificare
l'intero registro senza preoccuparci degli altri bit.

Non abbiamo bisogno di risolvere oggi un problema che ancora non esiste, ma è utile sapere che quella semplicissima riga
potrebbe non rimanere identica per tutta la vita del progetto.

Nessun ingrandimento

Il VIC-II permette anche di raddoppiare le dimensioni di uno sprite in orizzontale e in verticale. La nostra automobile,
per il momento, deve invece essere visualizzata nelle sue dimensioni normali, perciò inizializziamo esplicitamente
entrambi i registri:

310 rem nessun ingrandimento

320 poke 53271,0

330 poke 53277,0

53271 controlla l'espansione verticale degli sprite, mentre 53277 quella orizzontale. Anche in questo caso un singolo
bit corrisponde a ciascuno degli otto sprite e scrivere zero significa disabilitare l'espansione per tutti.

Potremmo apparentemente omettere queste istruzioni, confidando nel fatto che dopo l'accensione i registri si trovino già
nella condizione che desideriamo. Abbiamo preferito non farlo. Il nostro programma deve cercare di dichiarare
esplicitamente lo stato hardware dal quale vuole partire, evitando per quanto possibile di dipendere da ciò che è
accaduto prima della sua esecuzione.

È un'abitudine che diventerà sempre più importante man mano che il listato crescerà. Un programma facile da capire non è
soltanto un programma nel quale le istruzioni sono ordinate: è anche un programma che lascia meno stati impliciti
possibile.

Perché multicolor?

Se ci limitassimo alla modalità standard, ogni punto dello sprite potrebbe trovarsi in una delle due condizioni
possibili: trasparente oppure colorato. Avremmo quindi una buona risoluzione orizzontale, ma un solo colore visibile.

Per la nostra automobile abbiamo preferito una soluzione diversa.

350 rem sprite 0 in multicolor

360 poke 53276,1

Il registro 53276 stabilisce quali sprite devono utilizzare la modalità multicolor. Scrivendo 1 impostiamo il bit
relativo allo sprite 0 e lasciamo disattivati gli altri sette.

In modalità multicolor i bit vengono letti a coppie. Di conseguenza non abbiamo più 24 punti orizzontali indipendenti,
ma 12 elementi grafici più larghi, ciascuno rappresentato da due bit. In cambio possiamo attribuire alle quattro
combinazioni possibili significati differenti: trasparenza e tre colori visibili.

È un compromesso molto tipico della grafica del Commodore 64. Rinunciamo a una parte della risoluzione orizzontale per
ottenere maggiore ricchezza cromatica.

Nel nostro caso è una scelta sensata. L'automobile è vista dall'alto, ha dimensioni ridotte e deve essere immediatamente
riconoscibile sullo schermo. Disporre del rosso per la carrozzeria, del bianco e del nero per alcuni dettagli ci offre
un risultato più leggibile del semplice sprite monocromatico.

I tre colori vengono configurati con le istruzioni successive:

380 rem multicolor 1 = bianco

390 poke 53285,1

400 rem multicolor 2 = nero

410 poke 53286,0

420 rem colore sprite 0 = rosso

430 poke 53287,2

Gli indirizzi 53285 e 53286 contengono i due colori multicolor condivisi dagli sprite che utilizzano questa modalità.
Abbiamo scelto rispettivamente il bianco, codice 1, e il nero, codice 0.

Il registro 53287 contiene invece il colore specifico dello sprite 0, che impostiamo a 2, cioè rosso.

Questo significa che l'immagine memorizzata nei 63 byte non contiene direttamente i nomi o i codici dei colori. Contiene
soltanto coppie di bit. È il VIC-II, leggendo quei bit insieme ai registri che abbiamo appena configurato, a stabilire
quali punti devono essere trasparenti, quali bianchi, quali neri e quali rossi.

La distinzione è importante: i dati definiscono la forma; i registri contribuiscono a stabilire come quella forma deve
essere visualizzata.

Se domani decidessimo che l'automobile deve essere blu anziché rossa, non sarebbe necessario ridisegnare lo sprite.
Basterebbe modificare il valore scritto nel registro del suo colore.

Esiste già, ma non possiamo ancora vederla

A questo punto abbiamo fatto quasi tutto. I dati dell'automobile sono stati copiati nella memoria, il puntatore indica
il blocco corretto, abbiamo scelto posizione, dimensioni, modalità grafica e colori.

Eppure manca ancora un'istruzione:

450 rem abilita lo sprite 0

460 poke 53269,1

Il registro 53269 decide quali degli otto sprite devono essere visibili. Anche questa volta ogni bit corrisponde a uno
sprite; il valore 1 attiva il bit zero e quindi abilita lo sprite 0.

Soltanto quando questa istruzione viene eseguita il VIC-II inizia realmente a mostrarci ciò che abbiamo preparato.

È un dettaglio che mi piace particolarmente di questa prima milestone, perché rende molto chiara la distinzione tra
costruire qualcosa e renderlo visibile. L'automobile esisteva già nei 63 byte memorizzati a partire dall'indirizzo 832.
Il VIC-II conosceva già il blocco nel quale cercarla, la sua posizione e i suoi colori. Mancava soltanto il consenso a
visualizzarla.

Ora quel consenso è arrivato.

Sul monitor compare finalmente la nostra automobile.

Non corre ancora, non può essere guidata, non esistono avversari e intorno a lei non c'è nemmeno una strada. Se
giudicassimo la milestone dalla quantità di azione presente sullo schermo, il risultato potrebbe sembrare quasi
insignificante. Dal punto di vista del programma, però, abbiamo già attraversato parecchi concetti che ci
accompagneranno per tutto il progetto.

I numeri che disegnano la macchina

Rimane da vedere ciò che abbiamo lasciato deliberatamente alla fine del listato: i dati grafici veri e propri.

Nella M01 sono organizzati così:

480 rem ===============================

490 rem dati grafici dello sprite

500 rem 21 righe x 3 byte = 63 byte

510 rem ===============================

520 data \ \ 0, \ \ 0, \ \ 0, \ \ 0, \ \ 0, \ \ 0

530 data \ \ 3, 255, 192, \ \ 0, 170, \ \ 0

540 data \ \ 0, \ 40, \ \ 0, \ \ 1, 170, \ 64

550 data \ \ 3, 150, 192, \ \ 3, 150, 192

560 data \ \ 3, 150, 192, \ \ 0, 150, \ \ 0

570 data \ \ 0, 150, \ \ 0, \ \ 0, 170, \ \ 0

580 data \ \ 0, 190, \ \ 0, \ \ 5, 186, \ 80

590 data \ 15, 190, 240, \ 15, 170, 240

600 data \ 15, 255, 240, \ \ 2, 170, 128

610 data \ \ 2, 170, 128, \ \ 0, \ \ 0, \ \ 0

620 data \ \ 0, \ \ 0, \ \ 0

A prima vista questa è probabilmente la parte meno leggibile dell'intero programma. Una sequenza come 15,190,240 non
assomiglia affatto a un pezzo di automobile.

Eppure è proprio questo il punto nel quale la rappresentazione grafica diventa memoria.

Ogni gruppo di tre valori descrive una riga dello sprite. Ciascun valore decimale rappresenta un byte e quindi otto bit;
tre byte consecutivi forniscono i 24 bit necessari per una riga completa. In modalità multicolor quei bit vengono poi
interpretati a coppie, producendo i dodici elementi orizzontali che possono assumere uno dei quattro significati
cromatici disponibili.

Non è necessario imparare a riconoscere a colpo d'occhio una carrozzeria dentro la sequenza 15,190,240. Quello che ci
interessa comprendere è che l'automobile visualizzata sul monitor non esiste, per il computer, come concetto. Il
Commodore non conosce una ruota, un parabrezza o un cofano: conosce soltanto numeri che, trasformati in bit e
interpretati dal VIC-II secondo le regole che abbiamo configurato, producono quell'immagine.

Questa distanza tra ciò che vediamo noi e ciò che vede la macchina sarà uno dei temi ricorrenti del nostro lavoro.

Noi vediamo un'automobile.

Il Commodore vede 63 byte.

Entrambe le descrizioni sono corrette; appartengono semplicemente a due livelli differenti dello stesso problema.

C'è anche una scelta di formattazione che vale la pena osservare. I valori sono stati allineati inserendo spazi dopo le
virgole, in modo che numeri di una, due o tre cifre risultino più facilmente confrontabili. Il BASIC non ne ha bisogno
e li ignorerebbe anche se fossero scritti tutti attaccati. Ne abbiamo bisogno noi.

Quando un listato deve essere letto, controllato e magari digitato manualmente, la disposizione del codice non è un
ornamento. È parte della sua qualità.

La M01 completa

Adesso possiamo finalmente guardare il programma per intero. Ogni sua parte dovrebbe risultare molto meno misteriosa di
quanto sarebbe stata se avessimo incontrato subito il listato completo.

10 rem ================================

20 rem m01 - compare l'automobile

30 rem sprite multicolor giocatore

40 rem ================================

50 rem

60 rem pulisce lo schermo

70 print chr\$(147)

80 rem bordo nero

90 poke 53280,0

100 rem sfondo grigio chiaro

110 poke 53281,15

120 rem

130 rem carica i 63 byte dello sprite

140 rem a partire da \$0340 = 832

150 for i=0 to 62

160 read a

170 poke 832+i,a

180 next i

190 rem

200 rem blocco sprite: 832/64 = 13

210 rem puntatore sprite 0 = 2040

220 poke 2040,13

230 rem

240 rem posizione iniziale

250 poke 53248,172

260 poke 53249,120

270 rem

280 rem azzera il nono bit della x

290 poke 53264,0

300 rem

310 rem nessun ingrandimento

320 poke 53271,0

330 poke 53277,0

340 rem

350 rem sprite 0 in multicolor

360 poke 53276,1

370 rem

380 rem multicolor 1 = bianco

390 poke 53285,1

400 rem multicolor 2 = nero

410 poke 53286,0

420 rem colore sprite 0 = rosso

430 poke 53287,2

440 rem

450 rem abilita lo sprite 0

460 poke 53269,1

470 rem

480 rem ===============================

490 rem dati grafici dello sprite

500 rem 21 righe x 3 byte = 63 byte

510 rem ===============================

520 data \ \ 0, \ \ 0, \ \ 0, \ \ 0, \ \ 0, \ \ 0

530 data \ \ 3, 255, 192, \ \ 0, 170, \ \ 0

540 data \ \ 0, \ 40, \ \ 0, \ \ 1, 170, \ 64

550 data \ \ 3, 150, 192, \ \ 3, 150, 192

560 data \ \ 3, 150, 192, \ \ 0, 150, \ \ 0

570 data \ \ 0, 150, \ \ 0, \ \ 0, 170, \ \ 0

580 data \ \ 0, 190, \ \ 0, \ \ 5, 186, \ 80

590 data \ 15, 190, 240, \ 15, 170, 240

600 data \ 15, 255, 240, \ \ 2, 170, 128

610 data \ \ 2, 170, 128, \ \ 0, \ \ 0, \ \ 0

620 data \ \ 0, \ \ 0, \ \ 0

Il listato contiene molti commenti rispetto alla quantità di istruzioni realmente eseguite. È una scelta intenzionale.
Le righe REM non servono al Commodore per capire il programma; servono a noi per ricordare perché una determinata
istruzione è stata scritta e per rendere riconoscibili le diverse sezioni anche quando il codice comincerà a crescere.

Questa attenzione alla leggibilità può sembrare eccessiva adesso che il programma occupa poche decine di righe. Sarebbe
invece molto più difficile introdurla quando il gioco ne conterrà centinaia. Stiamo quindi cercando di costruire fin
dall'inizio non soltanto un programma che funzioni, ma un programma che possiamo continuare a comprendere mentre
evolve.

Vale la pena osservare anche un'altra caratteristica dei POKE utilizzati in questa prima milestone. Alcuni registri,
come quelli relativi al nono bit della coordinata X, all'espansione, alla modalità multicolor e all'abilitazione,
contengono contemporaneamente un bit per ciascuno degli otto sprite. Oggi possiamo scrivere tranquillamente valori come
0 oppure 1 perché esiste un solo sprite e nessun'altra informazione deve essere preservata. Quando entreranno in scena
i veicoli avversari, la stessa strategia potrebbe cancellare involontariamente la configurazione degli altri sprite.

Non dobbiamo ancora risolvere questo problema.

Ci basta averlo visto arrivare.

È un buon esempio del modo in cui procederà il nostro progetto: una soluzione perfettamente valida in una milestone può
mostrare i propri limiti quando il programma cresce. Non significa necessariamente che fosse una cattiva soluzione.
Significa che il problema è cambiato e che dovremo imparare qualcosa di nuovo per affrontarlo.

È davvero già un videogioco?

Se guardiamo lo schermo dopo aver digitato RUN, vediamo soltanto una piccola automobile rossa, immobile su uno sfondo
grigio chiaro incorniciato di nero. Sarebbe difficile convincere qualcuno che questo sia già uno shoot'em up
automobilistico.

Eppure, per me, quel momento possiede un significato particolare.

Quando ero ragazzo avevo immaginato molte volte di realizzare un videogioco sul Commodore 64, ma non avevo mai scritto
la prima riga di quel progetto. Questa volta, invece, sullo schermo c'è qualcosa che non appartiene più soltanto
all'immaginazione. È piccolo, non fa ancora nulla e non assomiglia nemmeno lontanamente al gioco che vogliamo
costruire, ma esiste perché abbiamo cominciato a trasformare un'idea in istruzioni.

Il professor Rossi Brunori, probabilmente, avrebbe trovato qualche modo molto più semplice per dirlo. Forse avrebbe
guardato quella macchina immobile e ci avrebbe chiesto che cosa mancasse perché diventasse davvero interessante. La
risposta, questa volta, sarebbe arrivata facilmente.

Deve muoversi.

Ed è esattamente il problema che affronteremo nella prossima milestone.


\bigskip
\end{document}
